% !PRINTED 199
\weberpage{199}
%<<CONT
\gesperrt{cienten $M$ dargestellt werden durch einen der beiden
Ausdrücke}
\[
\Sigma \, M D^\gamma H^\alpha f^\beta, \quad Q \, \Sigma \, M D^\gamma H^\alpha f^\beta.
\]
Die Invarianten sind als Specialfall unter den Covarianten
enthalten und sind sämmtlich Potenzen von $D$.

\begin{center}
\S.\,64.

\gesperrt{Biquadratische Formen.}
\end{center}

Wir gehen nun zur Betrachtung der biquadratischen Form
\weq{1}{f(x, y) = a_0 x^4 + a_1 x^3 y + a_2 x^2 y^2 + a_3 x y^3 + a_4 y^4}
über.

Wir haben zunächst die \gesperrt{Hesse}'sche Determinante als Covariante vierter Ordnung, die wir, um sie von einem Zahlenfactor
zu befreien, durch 3 theilen:
\weq{2}{H = \frac{1}{3}
\begin{vmatrix}
12 a_0 x^2 + 6 a_1 x y + 2 a_2 y^2, & 3 a_1 x^2 + 4 a_2 x y + 3 a_3 y^2 \\
3 a_1 x^2 + 4 a_2 x y + 3 a_3 y^2, & 2 a_2 x^2 + 6 a_3 x y + 12 a_4 y^2
\end{vmatrix},}
oder
\[
H = A_0 x^4 + A_1 x^3 y + A_2 x^2 y^2 + A_3 x y^3 + A_4 y^4,
\]
worin
\[
\begin{aligned}
A_0 &= 8 a_0 a_2 - 3 a_1^2, & A_4 &= 8 a_2 a_4 - 3 a_3^2, \\
A_1 &= 24 a_0 a_3 - 4 a_1 a_2, & A_3 &= 24 a_1 a_4 - 4 a_2 a_3, \\
A_2 &= 48 a_0 a_4 + 6 a_1 a_3 - 4 a_2^2, & &
\end{aligned}
\]
und eine Covariante sechsten Grades
\weq{3}{T = \frac{1}{12}
\begin{vmatrix}
f'(x), & f'(y) \\
H'(x), & H'(y)
\end{vmatrix}.}

Die Invarianten der biquadratischen Form bilden wir am
einfachsten nach \S.\,61 aus den Wurzeldifferenzen.

Wir erhalten so eine Invariante von der zweiten und eine
von der dritten Ordnung in den Coëfficienten:
\weq{4}{A = \tfrac{1}{2} a_0^2 [(12)^2(34)^2 + (13)^2(24)^2 + (14)^2(23)^2],}
\weq{5}{B = a_0^3 \, \Sigma \, (12)^2(34)^2(13)(42),}
Ausdrücke, die sich übersichtlicher schreiben lassen, wenn man
\weq{6}{U = (12)(34), \quad V = (13)(42), \quad W = (14)(23)}
setzt, nämlich
%CONT>>
